\label{chap:taxonomia-monodromia}

La representación posicional de un número registra una sucesión de símbolos,
pero no conserva necesariamente el estado que determina su continuación. La
Holografía Modular Triádica distingue ambos niveles. La fase nonádica indica
la posición de una lectura dentro del ciclo de nueve pasos; el estado
enriquecido conserva, además, orientación, cociente euclídeo, transporte de
arrastre y condición de frontera. El retorno de la fase no implica, por
tanto, el retorno del estado.

Sea \(S\) un espacio de estados enriquecidos, sea
\(T:S\to S\) una transición y sea
\[
g:S\longrightarrow\mathbb Z/9\mathbb Z
\]
la fase, con \(g(Ts)=g(s)+1\). Para una base \(b\geq2\), un lector
\(d:S\to\{0,\ldots,b-1\}\) produce
\[
x(s)=\sum_{n\geq0}\frac{d(T^ns)}{b^{n+1}}.
\]
Cuando un valor admite las dos expansiones posicionales usuales, se adopta la
representación canónica que no termina en una cola constante \(b-1\). Esta
convención separa una ambigüedad de escritura de la multiplicidad propiamente
HMT, que procede de secciones enriquecidas distintas con una misma
evaluación.
El operador de retorno nonádico es la restricción
\[
M=T^9:g^{-1}(r)\longrightarrow g^{-1}(r).
\]
La diferencia entre \(T^9s=s\) y \(g(T^9s)=g(s)\) formaliza la diferencia
entre periodicidad y monodromía.

La transición del estado enriquecido de TPK puede estar dada por una relación y no por una
función. En tal caso, la notación anterior se aplica al desplazamiento sobre
el espacio de historias infinitas o a una realización determinista cuyo estado
se haya ampliado con la información necesaria para fijar la continuación; no se
presupone que un bloque visible posea un sucesor único.

\begin{definition}[Número HMT y presentación aritmética]
Un número HMT es únicamente una sección compatible
\(\boldsymbol x\) del sistema inverso de estados supervivientes. Elegido un
lector posicional \(L\) cuyo dominio contiene \(\boldsymbol x\), el par
\((\boldsymbol x,L)\) es una presentación o medición del número. El valor
arquimediano publicado por esa presentación es la intersección puntual de los
cilindros asociados cuando sus diámetros tienden a cero. Dos secciones con el
mismo valor no se identifican antes de aplicar el cociente arquimediano.
\end{definition}

\begin{definition}[Periodicidades visible y enriquecida]
La salida es eventualmente periódica cuando existen \(N\geq0\) y \(p>0\)
tales que
\[
d(T^{n+p}s)=d(T^ns),\qquad n\geq N.
\]
El estado es eventualmente periódico cuando
\[
T^{N+p}s=T^Ns.
\]
La segunda condición implica la primera. La recíproca exige que el lector
separe las colas de la órbita considerada, es decir, que
\[
 \bigl(d(T^{n+k}s)\bigr)_{k\geq0}
 =\bigl(d(T^{m+k}s)\bigr)_{k\geq0}
 \quad\Longrightarrow\quad T^ns=T^ms.
\]
\end{definition}

\begin{theorem}[Clasificación por retorno de fase]
\label{thm:clasificacion-retorno-fase}
Sea \(s\in S\).
\begin{enumerate}
  \item Si la órbita de \(s\) alcanza un estado cuya salida futura es
  constantemente nula, entonces \(x(s)\) posee una expansión terminante y es
  racional.
  \item Si el estado es eventualmente periódico, entonces la salida es
  eventualmente periódica y \(x(s)\in\mathbb Q\).
  \item Si el lector separa las colas de la órbita y la órbita de \(s\) no es
  eventualmente periódica, entonces la salida no es eventualmente periódica
  y \(x(s)\notin\mathbb Q\).
  \item La no periodicidad no distingue entre irracionalidad algebraica y
  trascendencia. Esta distinción requiere un certificado polinómico o un
  teorema de trascendencia adicional.
\end{enumerate}
\end{theorem}

\begin{proof}
Los dos primeros apartados son la caracterización clásica de los racionales
mediante expansiones posicionales terminantes o eventualmente periódicas. En
el tercero, una eventual periodicidad de la salida produciría dos colas
iguales. La hipótesis de separación identificaría entonces los estados
correspondientes, contradiciendo la no periodicidad de la órbita. El cuarto
apartado sigue de la existencia de irracionales algebraicos y trascendentes
con expansiones no eventualmente periódicas.
\end{proof}

\begin{corollary}[Monodromía nonádica y aperiodicidad]
Supóngase que la fase retorna cada nueve pasos, que la órbita de retorno
\[
 \{M^ks:k\geq0\},\qquad M=T^9,
\]
no es eventualmente periódica y que el lector separa las colas. Entonces la
salida no es eventualmente periódica y el número producido es irracional. La
periodicidad de la fase convive en este caso con una evolución aperiódica del
estado enriquecido.
\end{corollary}

La hipótesis no puede sustituirse por una sola desigualdad \(M(s)\ne s\), ni
por un número finito de retornos distintos. Relativamente al estado inicial,
la transición y el lector declarados, la ventana vigente certifica nueve
retornos de fase con cambio de estado; demuestra monodromía finita, pero no
por sí sola la aperiodicidad de toda la órbita.

\ifdefined\HMTRepetirResultadosTaxonomicos
\section{El cuadrado local de \texorpdfstring{\(e\)}{e} y la ruptura de memoria}

Para una palabra decimal \(a_1\cdots a_n\), llamaremos \emph{cuadrado local}
de longitud \(\ell\) en la posición \(i\) a una factorización
\[
 a_i\cdots a_{i+2\ell-1}=vv,
 \qquad |v|=\ell.
\]
Esta noción pertenece a combinatoria de palabras y no debe confundirse con
un período de la cola infinita.

\begin{theorem}[Caracterización finita del cuadrado \(1828^2\)]
Considérense las \(243\) palabras del catálogo N33 y el lector congelado
\(w_6\mapsto w_{30}\mapsto d_1d_2d_3d_4\), con cuatro tríadas en base mil.
Existe una única palabra cuya lectura contiene un cuadrado de longitud cuatro
iniciado en la segunda cifra decimal. Es
\[
 w_e=201101,
 \qquad
 718281828459=7\,\underbrace{1828\,1828}_{(1828)^2}\,459.
\]
La repetición no se prolonga a una tercera copia: la cifra siguiente requerida
sería \(1\), mientras que la cifra efectiva es \(4\). En todo el catálogo sólo
existe otro cuadrado de longitud cuatro,
\[
 575157518472=\underbrace{5751\,5751}_{\text{cuadrado de palabra}}\,8472,
\]
pero comienza en la primera cifra.
\end{theorem}

\begin{proof}
La enumeración exhaustiva evalúa las \(243\) palabras mediante aritmética
entera exacta. El perfil de cuadrados para longitudes \(1,\ldots,6\) contiene,
respectivamente,
\[
240, 21, 3, 2, 0, 0
\]
ocurrencias, distribuidas entre
\[
153, 19, 3, 2, 0, 0
\]
palabras. Las dos ocurrencias de longitud cuatro son exactamente las escritas
en el enunciado. Dos enumeraciones enteras independientes reproducen el mismo
censo y sus salidas completas coinciden término a término.
\end{proof}

La palabra \(w_e\) posee además la cadena de elevación
\[
201101\mid121221\mid102011\mid012222\mid102011.
\]
El tercer y el quinto bloque visibles coinciden. No coinciden, sin embargo,
los estados que los transportan. Antes de ambas emisiones, sus reducciones
módulo \(27\) son
\[
(25,11,7,17,8,16),
\qquad
(7,8,4,23,26,7),
\]
y los cocientes de Hensel posteriores son
\[
(6,13,9,2,13,1),
\qquad
(15,0,3,19,23,25).
\]
Por tanto, el retorno \(102011\) es un retorno de la proyección visible, no
del estado enriquecido. Esta es la formulación exacta de la intuición de una
repetición local que se rompe: el cociente recuerda una genealogía que el
bloque residual no conserva.

El resultado pertenece al lector enriquecido ya fijado por el sistema
APP--TRIT--TPK\@. El censo finito caracteriza exhaustivamente la aparición del
cuadrado local; la prolongación coinductiva de la sección pertenece a la
conexión nonádica construida en el
\cref{chap:numero-heisenberg-d108}. La propiedad combinatoria del prefijo y
la generación de la expansión completa son, por tanto, dos niveles
compatibles de una misma genealogía.

Esta formulación proporciona el marco exacto para interpretar el comienzo
decimal de \(e\),
\[
e=2.718281828459\ldots.
\]
La expresión correcta ya no es ``semiperiodicidad'', sino \emph{cuadrado local
de palabra con retorno visible y ruptura de memoria}.

\section{2 medidas finitas del emisor \(468\to243\): frecuencia y prioridad estructural}

No existe una distribución uniforme sobre todos los números reales. Sí
existe una medida canónica sobre el catálogo finito cuando se toma uniforme
el conjunto de \(104\,976\) semillas TPK\@. Para una palabra visible \(w\), se
define
\[
\mu_{\mathrm{TPK}}(w)
=\frac{\#\{\text{semillas que proyectan sobre }w\}}{104\,976}.
\]
La medida no es uniforme. Las \(243\) palabras poseen pesos comprendidos
entre \(144\) y \(2088\), con nueve valores de multiplicidad distintos.
El histograma exacto, escribiendo el peso como base y el número de palabras
como exponente, es
\[
144^{120},\ 288^{54},\ 432^{18},\ 576^9,\ 864^{15},\
1008^6,\ 1440^3,\ 1944^{12},\ 2088^6.
\]

Para las tres palabras distinguidas se obtiene
\[
\mu_{\mathrm{TPK}}(010211)=\frac{1008}{104\,976}=\frac7{729},
\]
\[
\mu_{\mathrm{TPK}}(201101)
=\mu_{\mathrm{TPK}}(121200)
=\frac{144}{104\,976}=\frac1{729}.
\]
La palabra asociada a \(\pi\) tiene cinco preimágenes \(U_6\) y peso siete
veces mayor que las palabras asociadas a \(e\) y \(\varphi\). Este resultado
demuestra una jerarquía combinatoria en el catálogo; no identifica por sí
solo las palabras con las constantes clásicas ni establece una probabilidad
sobre \(\mathbb R\).

Tampoco define por sí solo una prioridad total: seis palabras tienen peso
\(1008\), mientras que \(e\) y \(\varphi\) pertenecen a la clase mínima de
peso \(144\), compartida por otras \(118\) palabras. La distinción HMT debe
expresarse mediante un perfil estructural, por ejemplo
\[
\mathcal P(w)=
\bigl(m_{\rm semilla}(w),m_{U_6}(w),r_{\rm supervivencia}(w),
s_{\rm frontera}(w),s_{\rm excepcional}(w)\bigr),
\]
con cada coordenada definida antes de aplicar etiquetas externas. Un escalar
de ``prioridad'' exige además fijar, también de antemano, cómo se ordenan o
ponderan esas coordenadas.
\fi

\section{Clasificación de las constantes por antecedente genealógico, operación generadora y codominio}

La clasificación HMT separa propiedades que la representación ordinaria
reúne bajo un solo valor:
\begin{center}
\begin{tabular}{p{0.25\textwidth}p{0.64\textwidth}}
\toprule
Clase & Condición estructural \\
\midrule
Terminante & La salida alcanza una cola nula. \\
Racional periódico & La salida es eventualmente periódica. \\
Racional monodrómico & La salida es periódica, pero el estado enriquecido
recorre una fibra no trivial que el lector visible olvida. \\
Irracional coinductivo & La salida no es eventualmente periódica y los
cilindros determinan un único valor. \\
Recurrente local & Una palabra finita contiene retornos o cuadrados, sin que
ello imponga periodicidad a la cola. \\
Algebraico & El valor satisface un polinomio no nulo con coeficientes enteros;
la ruta debe acompañarse de ese certificado. \\
Trascendente & El valor no satisface ningún polinomio entero no nulo; la
clasificación exige un teorema de trascendencia. \\
Distinguido HMT & La sección satisface un predicado estructural congelado de
fibra, frontera, supervivencia o simetría; es una clase transversal a las
anteriores. \\
\bottomrule
\end{tabular}
\end{center}

Así, HMT no altera las definiciones aritméticas correctas de racionalidad,
algebraicidad y trascendencia. Las refina añadiendo la genealogía del estado,
la monodromía de fase y la multiplicidad de las fibras que la evaluación
arquimediana había olvidado.

En particular, \(\varphi\) es algebraico porque satisface
\(X^2-X-1=0\), mientras que la trascendencia de \(e\) y \(\pi\) procede de
teoremas externos clásicos. Estas propiedades son compatibles con que las tres
secciones sean distinguidas por HMT\@. El cierre dodecafásico construye más
adelante la proyección racional exacta
\(\alpha_{12}=\pi_{12}(\alpha_{\mathrm{HMT}})\), y la prolongación compatible
construye la sección coinductiva \(\alpha_{\mathrm{HMT}}\), en el
\cref{chap:alpha-rev7}; ninguna de ellas se usa en esta taxonomía como entrada
ni se publica aquí como una segunda residencia. El núcleo \(E_{21/22}\)
produce una raíz analítica distinta, cuya concordancia con la vía entera se
demuestra en el codominio de \(\operatorname{Tr}_9\). Sección, proyección
racional, raíz analítica, lector y valor conservan así sus tipos propios.

Finalmente, una ``dirección infinitesimal'' requiere tipado propio. Para un
lector finito \(r_m\), dos secciones distintas en una misma fibra de \(r_m\)
definen una diferencia invisible a profundidad \(m\), aunque activa en una
prolongación posterior. No constituyen por ello un real infinitesimal no nulo.
Si dos secciones coinciden bajo todas las truncaciones fieles del sistema
inverso, son la misma sección. Un infinitesimal aritmético genuino exigiría
construir una extensión ordenada o valuada no arquimediana y demostrar que las
operaciones HMT descienden a ella.
